\documentclass[10pt,a4paper]{amsart}
\usepackage[margin=19mm]{geometry}
\usepackage{amsmath,amssymb,mathtools}
\usepackage[scr=rsfs,cal=euler]{mathalfa}
\usepackage{xfrac}
\usepackage[unicode,hidelinks,bookmarks=false]{hyperref}
\newcommand{\ie}{i.e.\ }
\newcommand{\cf}{cf.\ }
\newcommand{\resp}{respectively\ }
\newcommand{\Subsection}{\subsection}
\providecommand{\nobreakdash}{\nobreak\hbox{-}}
\setlength{\emergencystretch}{2em}
\multlinegap0pt
\usepackage{bm}
\usepackage{bbm}
\usepackage{dsfont}
\usepackage{xspace}
\usepackage{cancel}
\usepackage{mathtools}
\usepackage{amsthm}
\makeatletter
\@ifundefined{theorem}{\newtheorem{theorem}{Theorem}}{}
\@ifundefined{lemma}{\newtheorem{lemma}[theorem]{Lemma}}{}
\@ifundefined{proposition}{\newtheorem{proposition}[theorem]{Proposition}}{}
\makeatother
\let\mathcal\mathscr

\numberwithin{equation}{section}
\newcommand{\NN}{\mathbb{N}}
\newcommand{\bgamma}{\boldsymbol{\gamma}}
\newcommand{\btheta}{\boldsymbol{\theta}}
\renewcommand{\a}{\mathbf{a}}
\renewcommand{\c}{\mathbf{c}}
\renewcommand{\leq}{\leqslant}
\title[Integral solutions to systems of diagonal equations]{Integral solutions to systems of diagonal equations}
\author{Nick Rome, Shuntaro Yamagishi}
\date{Source: arXiv:2406.09256v2 (CC BY 4.0)}
\begin{document}
\maketitle

\noindent\emph{Source and licence.} Integral solutions to systems of diagonal equations. Version arXiv:2406.09256v2 (CC BY 4.0), distributed under the Creative Commons Attribution 4.0 International licence, as recorded in the arXiv licence metadata for this exact version and shown on the abstract page https://arxiv.org/abs/2406.09256v2. Full rights notice: licenses/arxiv-2406.09256-CC-BY-4.0.txt.\par\smallskip

\noindent\textit{Editorial scope.} Counted unit: the Proposition whose source label is \texttt{lem:major} in the article (source0.tex lines 930--938, source label \texttt{lem:major}), delivered together with the article's own complete original proof of it (source0.tex lines 940--995, one \texttt{proof} environment of 56 lines). No mathematics is written, repaired or re-derived: statement and proof are the article's own text line for line. Required context retained uncounted: lem33 (lines 898--918); CB (lines 901--907). Every definition, display and label of those lines is retained unchanged. Omitted from this package and not redistributed: 1--897, 908--929, 939--939, 996--1308. Nothing inside a retained excerpt is omitted or altered: every retained body is delivered exactly as the article's lines are written, so no omission receipt of any kind accompanies this package. Citations: the retained excerpts carry no citation command at all, so no bibliography entry of the article is transcribed and no reference is redistributed. Reference adaptation: none was needed and none was made; every reference inside the delivered unit resolves inside it, so no unresolved reference or citation is produced. Printed slips kept literal: no printed word, symbol, hypothesis or number of the counted statement or of its proof was altered. Layout and class: the article's own class is \texttt{amsart} with packages including amsfonts, amsthm, amsmath, amssymb, bbm, bm, color, enumerate, tikz-cd, eurosym, url, hyperref, xy, geometry; the wrapper is the lane's pinned \texttt{amsart} editorial wrapper at 10pt on A4, which restates the theorem-like environments the retained text needs and adds the article's own macro definitions that the retained text needs (\texttt{\textbackslash NN}, \texttt{\textbackslash a}, \texttt{\textbackslash bgamma}, \texttt{\textbackslash btheta}, \texttt{\textbackslash c}, \texttt{\textbackslash leq}, \texttt{\textbackslash mathcal}), with extra wrapper packages (bm, bbm, dsfont, xspace, cancel, mathtools). The delivered numbering is the unit's own. Assets: no retained span contains a figure environment, a graphics-inclusion command, a PDF-page inclusion, a table or a TikZ picture; no image file is copied into this package, the package declares no asset dependency and no image file is redistributed with it. Licence: version arXiv:2406.09256v2 of this article is distributed under the Creative Commons Attribution 4.0 International licence, as recorded in the arXiv licence metadata for this exact version and shown on the abstract page https://arxiv.org/abs/2406.09256v2. A full rights notice is delivered at licenses/arxiv-2406.09256-CC-BY-4.0.txt. Characters: every retained excerpt is pure ASCII, so the delivered engine, XeLaTeX with its default Latin Modern text font, reports no missing character.
\begin{lemma}
\label{lem33}
Let $\eta > 0$ be sufficiently small and
\begin{eqnarray}
\label{CB}
C_\mathfrak{B} = \begin{cases}
                   1 & \mbox{if } \mathfrak{B} = \NN, \\
                   \varrho(1/\eta) & \mbox{if } \mathfrak{B} = \mathcal{A}(X, X^{\eta}).
                 \end{cases}
\end{eqnarray}
Let $\delta, A > 0$ be sufficiently small.
Suppose that $0 \leq a \leq q \leq \mathfrak{L}$, $\beta = \alpha - a/q$ and $|\beta| < \mathfrak{L} X^{-d}$.
Then
$$\left \vert
\sum_{ x \in \mathfrak{B} \cap [1, X]}  e(\alpha x^d) - C_{\mathfrak{B}} X q^{-1} S(q,a) I(X^d \beta) \right \vert \ll
\begin{cases}
\mathfrak L^2 \phantom{pp}   &\mbox{if } \mathfrak B = \mathbb N,\\
\frac{X \mathfrak L^2}{\log X} &\mbox{if }\mathfrak B = \mathcal{A}(X, X^\eta).
\end{cases}
$$
\end{lemma}

\begin{proposition}\label{lem:major}
Let $\eta > 0$ be sufficiently small and $C_\mathfrak{B}$ as in (\ref{CB}).
Then
\[
\int_{\mathfrak{M}_\mathfrak{L} ^+} \prod_{\c \in \textnormal{Col}(M)} S_\c(\btheta) \cdot  e\left(-\sum_{i = 1}^{R}\mu_i \theta_i  \right)\mathrm{d}\btheta
=
C_\mathfrak{B}^n  X^{n - d R} \mathfrak{S}( \mathfrak L) \mathfrak{I}(\mathfrak L) + O( X^{n  - d R} \mathfrak{L}^{- 1}).
\]
\end{proposition}

\begin{proof}
First, if $\btheta \in \mathfrak{M}_\mathfrak{L}^+$ then there exist $0 \leq \a \leq q$ such that $\gcd(q, \a) = 1$ and
$$
\left| \c \cdot \btheta - \frac{ \c \cdot \a }{q} \right|  < C \mathfrak{L} X^{-d},
$$
where $C > 0$ is a constant depending only on $\c$; therefore, $\c \cdot \btheta$, reduced modulo $1$, satisfies the hypotheses of  Lemma \ref{lem33} with
$\c \cdot \bgamma$ and $C \mathfrak{L}$ in place of $\beta$ and $\mathfrak{L}$, respectively.
Thus we may apply Lemma  \ref{lem33} to $S_\c(\btheta)$ for any $\c \in \textnormal{Col}( M )$ and $\btheta \in \mathfrak{M}_\mathfrak{L}^+$.
The measure of $\mathfrak M_{\mathfrak L}^+$ is at most $\mathfrak L^{2R+1} X^{-dR}$ and thus integrating the error term
coming from applying Lemma \ref{lem33} to
$\prod_{\c \in \textnormal{Col}( M )} S_\c(\btheta)$
gives a total error of size
\[
\begin{cases}
O(X^{n-dR-1} \mathfrak{L}^{2R+3})  &\mbox{if } \mathfrak B = \mathbb N,\\
O(X^{n-dR} \frac{\mathfrak{L}^{2R+3}}{\log X}) &\mbox{if } \mathfrak B = \mathcal A(X, X^\eta).
\end{cases}
\]
The former case clearly provides a suitable error term for sufficiently small $\delta >0$ and in the latter case the error suffices on choosing $A < \frac{1}{2R+4}$.
As a result we have
\begin{eqnarray}
\label{major}
\notag
&&
\int_{  \mathfrak{M}^+_\mathfrak{L}  }  \prod_{\c \in \textnormal{Col}( M )} S_\c(\btheta) \cdot  e \left(-\sum_{i=1}^{R}  \mu_i  \theta_i    \right) \mathrm{d}\btheta
\\
\notag
&=& C_{\mathfrak{B}}^n X^n \sum_{1 \leq q \leq \mathfrak{L}}  q^{- n }  \sum_{ \substack{ 1 \leq  \a \leq  q \\ \gcd(q, \a) = 1  } }
 \prod_{\c \in \textnormal{Col}( M )} S_\c (\a /q)  \cdot e_q \left(-\sum_{i=1}^{R}   \mu_i  a_i   \right) \times
\\
\notag
&&\int_{\vert \bgamma \vert <  \mathfrak{L} X^{-d}}
\prod_{\c \in \textnormal{Col}( M )} I ( X^d \c\cdot\bgamma )  \cdot   e \left(-\sum_{i=1}^{R}  \mu_i  \gamma_i     \right) \mathrm{d}\bgamma
+ O \left(  X^{n  - d R} \mathfrak{L}^{- 1}\right)\\
\notag
&=&
C_{\mathfrak{B}}^n  X^n \mathfrak{S} (\mathfrak{L}) \int_{\vert \bgamma \vert <  \mathfrak{L} X^{-d}}
\prod_{\c \in \textnormal{Col}( M )} I (X^d \c\cdot \bgamma )  \cdot   e \left(-\sum_{i=1}^{R}  \mu_i  \gamma_i    \right) \mathrm{d}\bgamma
+ O ( X^{ n - d R} \mathfrak{L}^{ -1 }  )
\\
\notag
&=&
C_{\mathfrak{B}}^n  X^{n- d R }  \mathfrak{S} (\mathfrak{L}) \int_{\vert \bgamma \vert <  \mathfrak{L} }
\prod_{\c \in \textnormal{Col}( M )} I( \c\cdot \bgamma )  \cdot   e \left( - \frac{ 1 }{X^d}   \sum_{i=1}^{R}  \mu_i \gamma_i    \right) \mathrm{d}\bgamma
+ O ( X^{ n - d R} \mathfrak{L}^{ -1 } )
\\
\notag
&=&
C_{\mathfrak{B}}^n  X^{n - d R} \mathfrak{S} (\mathfrak{L}) \mathfrak{I} (\mathfrak{L}) + O ( X^{ n - d R} \mathfrak{L}^{ -1 }),
\end{eqnarray}
which completes the claim.
%Lastly, we obtain from Lemmas \ref{SSSS} and \ref{singint} that
%$$
%\mathfrak{S} (\mathfrak{L}) \mathfrak{I} (\mathfrak{L})  =  \mathfrak{S}  \mathfrak{I}  +  O \left(  \mathfrak{L}^{ 1 - R (\frac{T}{d} - 1) } + \mathfrak{L}^{1 - \frac{T}{d} } \right).
%$$
\end{proof}

\end{document}
